% ###################
% ## Filename: sCM_drivers.tex
% ## Location: D:\Project\Prej_Teaching\ISE\SupplyChainMgt\presentation\scm_02_drivers
% ## Created: Dec 14, 2010
% ## Last updated: Dec 14, 2010
% ## Figure Location: ..\Figures\
% ## Note: Presentation, Teaching
% ###################


%\documentclass[10pt,presentation]{beamer}
%\documentclass[10pt,handout]{beamer}
\documentclass[10pt,handout,notes=show]{beamer}

%============================================================================================================--
\usepackage{fontspec}
\usepackage{xunicode}
\usepackage{xltxtra}
\usepackage{array}
\usepackage{xcolor}
\usepackage{verbatim}
\usepackage{beamerthemesplit}
\usepackage{beamerbasetemplates}
\usepackage{graphics}
\usepackage{graphicx}
\usepackage{hyperref}
\usepackage{multirow}
\usepackage{tikz,pgfplots}
\usepackage{algorithm}
\usepackage{algorithmic}
\usepackage{soul,xcolor}
\usepackage{colortbl}
\usepackage{multicol}

\XeTeXlinebreaklocale"th"
\defaultfontfeatures{Scale=1.5, Mapping=tex-text} %{Scale=1.23}



\usetheme{Madrid}  % JuanLesPins  Rochester
           %  Berkeley  Palo Alto    with sidebar and top
           % Goettingen  Marburg     with sidebar
           %  Copenhagen  Luebeck  Warsaw

\setbeamertemplate{headline}
{%
    \begin{beamercolorbox}{section in head/foot}
        \vskip-2pt \hskip5pt
        \insertsectionnavigationhorizontal{\textwidth}{}{}
        \hskip5pt \vskip-2pt
    \end{beamercolorbox}
}%
\setbeamertemplate{footline}
{%
    \begin{beamercolorbox}{section in head/foot}
        \vskip-2pt \hskip5pt
        \insertshorttitle \hfill \insertshortdate \hskip10pt \insertframenumber / \inserttotalframenumber
        \vskip-2pt\hskip5pt
    \end{beamercolorbox}
}%
\definecolor{alertedOrange}{RGB}{255,36,0}
\setbeamercolor{alerted text}{fg=alertedOrange}
\useinnertheme{rounded}
\usecolortheme{rose}
%\setbeamercovered{transparent}
\usefonttheme{structuresmallcapsserif}
\definecolor{LRed}{rgb} {1,.8,.8}
\definecolor{LBlue}{rgb}{.8,.8,1}

%%%%%%%%%%%%%%%%%%%%%%%%%%%%%%%%%%%%%%%%%%%%%%%%%%%%%%%%%%%%%%%%%%%%%%%%%%%%%%%
% Notation:
\newtheorem{observation}{Observation}
\newcommand{\figureLocation}[1]{../figures/#1}
\newcommand{\figureLoc}[1]{../figures/#1}
\newcommand{\refLocation}[1]{../../../../Reference/#1}
\newcommand{\NextLine}{\vskip 0.05\textheight}
\newcommand{\NextHalfLine}{\vskip 0.02\textheight}
\newcommand{\RemoveLine}{\vskip -0.1\textheight}
\newcommand{\RemoveHalfLine}{\vskip -0.02\textheight}
\newcommand{\BackTab}{\hskip -0.15\textwidth}
\newcommand{\HalfBackTab}{\hskip -0.11\textwidth}
\newcommand{\DualFigureSize}{1.10\textwidth}
\newcommand{\MiniPageSize}{0.52\textwidth}
\newcommand{\code}[1]{\textcolor{blue}{#1} }
\newcommand{\rightNote}[2]{\hfill \uncover<#1>{\textbf{\code{\#\# {#2}  }}}}
\newcolumntype{P}[1]{>{\raggedright\arraybackslash}p{#1}<{}}

\author[Kittithreerapronchai]{%
  Oran~Kittithreerapronchai\inst{1} }
\institute[IE Chula]{
  \inst{1}%
  Department of Industrial Engineering, Chulalongkorn University\\
  Bangkok 10330 THAILAND \\
  }
\subject{}
\date{}

\setbeamertemplate{enumerate items}[default]
\setbeamertemplate{bibliography item}[text]
\bibliographystyle{alpha}
\newcommand{\textFont}[1]{\footnotesize{#1}}
\setcounter{MaxMatrixCols}{20}

%%%--------------------------------

\makeatletter
\defbeamertemplate{note page}{infolines}
{%
  {%
    \small
    \insertvrule{.25\paperheight}{white!90!black}
    \vskip-.25\paperheight
    \nointerlineskip
    \vbox{
      \hfill\insertslideintonotes{0.25}\hskip-\Gm@rmargin\hskip0pt%
      \vskip-0.25\paperheight%
      \nointerlineskip
      \begin{pgfpicture}{0cm}{0cm}{0cm}{0cm}
        \begin{pgflowlevelscope}{\pgftransformrotate{90}}
          {\pgftransformshift{\pgfpoint{-2cm}{0.2cm}}%
          \pgftext[base,left]{\footnotesize\the\year-\ifnum\month<10\relax0\fi\the\month-\ifnum\day<10\relax0\fi\the\day}}
        \end{pgflowlevelscope}
      \end{pgfpicture}}
    \nointerlineskip
    \vbox to .25\paperheight{\vskip0.5em
      \hbox{\insertshorttitle[width=8cm]}%
      \setbox\beamer@tempbox=\hbox{\insertsection}%
      \hbox{\ifdim\wd\beamer@tempbox>1pt{\hskip4pt\raise3pt\hbox{\vrule
            width0.4pt height7pt\vrule width 9pt
            height0.4pt}}\hskip1pt\hbox{\begin{minipage}[t]{7.5cm}\def\breakhere{}\insertsection\end{minipage}}\fi%
      }%
      \setbox\beamer@tempbox=\hbox{\insertsubsection}%
      \hbox{\ifdim\wd\beamer@tempbox>1pt{\hskip17.4pt\raise3pt\hbox{\vrule
            width0.4pt height7pt\vrule width 9pt
            height0.4pt}}\hskip1pt\hbox{\begin{minipage}[t]{7.5cm}\def\breakhere{}\insertsubsection\end{minipage}}\fi%
      }%
      \setbox\beamer@tempbox=\hbox{\insertshortframetitle}%
      \hbox{\ifdim\wd\beamer@tempbox>1pt{\hskip30.8pt\raise3pt\hbox{\vrule
            width0.4pt height7pt\vrule width 9pt
            height0.4pt}}\hskip1pt\hbox{\insertshortframetitle[width=7cm]}\fi%
      }%
      \vfil}%
  }%
  \vskip.25em
  \nointerlineskip
  %\setmainfont{TH Sarabun New}
  \begin{minipage}{\textwidth} % this is an addition
  \footnotesize{\fontspec{TH Sarabun New} \insertnote}
  \end{minipage}               % this is an addition
  %\setmainfont{Computer Modern}

}
\makeatother

\setbeamertemplate{note page}[infolines]

\title[LSCM v3.0]{Lecture 03:\\ Decision and Drivers in Supply Chain Management}
\hypersetup{
  pdftitle={03 Decision and Drivers in Supply Chain Management},
  pdfsubject={2104509lscm & 2184408scm handout},
  pdfproducer={XeLatex/WinEdt},
  pdfcreator={oran.k@chula.ac.th}
}


\begin{document}

%==============================================================================================================
\begin{frame}[t]
  \titlepage \label{Slide:Title}


  \NextLine   \NextLine
  \begin{flushright}
    \tiny{ \textbf{last updated:} \today } \hskip3ex
  \end{flushright}

  \note{
    \textbf{Purpose:} Transition from SCM strategy to SCM operation   \\
    \textbf{Leading Question:} How do we keep what we learn in the previous class and apply thme? \\
    \begin{itemize}
      \item \textbf{BBA:} For BBA (non-engineering) student, this is their SCM last topic
      \item \textbf{Decision:} How do we view/formulate SCM decison/problem? What are scopes?  \textbf{hint:} same with OM
      \item \textbf{Drivers:} How key SCM costs (facility, inventory, transportation) and response time affect by \# stores? What are other factors invovled?
      \item \textbf{Network:} Development of distribution network.
    \end{itemize}
  }
\end{frame}

%==============================================================================================================
\begin{frame}[t]
  \frametitle{Outline}

  \tableofcontents

  \vfill
  \begin{flushright}
    \tiny{\textbf{General Reference:} \cite{jacobs2010operations}~\cite{ballou2007evolution} \cite{chopra2007supply} \cite{goetschalckx2011supply} } \hskip3ex
  \end{flushright}

  \note{
    \textbf{Purpose:} Outline topics\\
    \textbf{Leading Question:} $<<$none$>>$\\
    \begin{itemize}
      \item \textbf{Decision:} how IE/OM from topic and make SCM decision. What are popular decisions that you may face?
      \item \textbf{7Eleven:} an example of case study:  (from US perspective, this is radical retailer)
      \item \textbf{Conceptual Design:} In general, what are factors that you should think before design network? and how they interact?
      \item \textbf{Conceptual Distribution:} how 7Eleven JPN distribute? and general development
    \end{itemize}
  }
\end{frame}

%=================================================================================================
\section[Decisions]{Driven Factors and Decisions in Supply Chain Management}
%============================================================================================================--
\begin{frame}
  \frametitle{Supply chain as Math modeling}


  \begin{itemize}
    \item \textbf{Stakeholder:} retailer, shipper, \alert{consumer}
    \item \textbf{Objective:} maximize \alert{value} of \alert{every stakeholder}
    \item \textbf{Decision:}
    \begin{itemize}
      \item \textbf{Supply Mgt:} \alert{sourcing}, strategic alliance, inventory
      \item \textbf{Product Mgt:} make or buy, MRP, JIT
      \item \textbf{Transportation Mgt:} location, mode-vehicle, routing
      \item \textbf{Integration/Collaboration:} IT (bullwhip effect), EDI, virtual integration \\[5ex]
    \end{itemize}
    \item \textbf{Information:}  forecasting, supplier evaluation, sales, network
    \item \textbf{Constraints:} \alert{law}, technology-infrastructure, \alert{time-capacity}, service level, quality
  \end{itemize}


  \note{
    \textbf{Purpose:} Systematic view of SCM decision\\
    \textbf{Leading Question:} Why do an OM/OR person specialize in SCM? \\
    \begin{itemize}
      \item \textbf{B/G} OR as Science of Better
      \item \textbf{Stakeholder:} how to form problem/desion? do we have to care? Which shoes/roles you are in?
      \item \textbf{Objective:} how to objectively measure \alert{value}. Utility $\gg$ Profit $\gg$ Cost
      \item \textbf{Decision:} four related topics
      \item \textbf{Information:} integrated through information, planning, and integration activities. how to obtain? does it reliable?
      \item \textbf{Constraint:} what we should aware?
      \begin{itemize} \scriptsize{
        \item \textbf{Law:} load/weight, zoning, speed, load checking
        \item \textbf{Tech:} lift, hoist, GPS, inventory, RFID, software
        \item \textbf{Time:} lead time, cycle time, seasonal,
        \item \textbf{misc:} capacity \& service lv \& quality \& quantity }
      \end{itemize}
      \item \textbf{Cheating:} usually we need to focus $\rightarrow$ which one?
    \end{itemize}
  }
\end{frame}

%==============================================================================================================
\begin{frame}
  \frametitle{Decision/Driver in Supply Chain}

  \textsc{Time Horizon}
  \begin{itemize}
    \item \textbf{SC Strategy:} structure of the supply chain; \alert{long-term}
    \item \textbf{SC Planning:} policies governing short-term operations; \alert{few months}
    \item \textbf{SC Operation:} \alert{weekly} to maximize efficiency of supply chain
  \end{itemize}

  \NextLine
  \textsc{Drivers in Supply Chain}
  \begin{itemize}
    \item \textbf{Facilities:} \alert{location} for manufacturing products and housing inventory
    \item \textbf{Inventory:} raw materials, WIP, F/G to \alert{smooth operations}
    \item \textbf{Transportation:} \alert{moving} products from $A$ to $B$
    \item \textbf{Information:} data and analysis to \alert{coordinate} between stages (best drivers)
    \item \textbf{Sourcing:} managing raw material and relationship with supplier
    \item \textbf{Pricing:} pricing goods and services provided by a firm to the supply chain
  \end{itemize}

  \note{
    \textbf{Purpose:} How to simplify decision making in SCM\\
    \textbf{Leading Question:} Do we have priority/cheat code
    \begin{itemize}
      \item \textbf{Supply Chain Strategy:} structure of the supply chain and locations; \alert{long-term} and \alert{expensive};
      \item[] \hskip3exLocations and capacities of facilities; Manufacturing and Storage locations; Modes of transportation \& IT systems \NextHalfLine
      \item \textbf{Supply Chain Planning:} tactical policies that govern short-term operations usually \alert{few months} or \alert{quarter}
      \item[] \hskip3ex Markets allocations and promotions; Inventory policies, Batch  sizing; Subcontracting, backup locations \NextHalfLine
      \item \textbf{Supply Chain Operation:}  \alert{weekly} or daily operations to maximize efficiency of supply chain
      \item[] \hskip3ex allocate orders to inventory, production, shipment; Set order due dates, routing, or delivery schedules;Replenish and product orders
    \end{itemize}
    \textbf{Driver:} inter-related topics in SCM (\textbf{note} continue next page)
    \begin{itemize}
      \item \textbf{Facilities:} manufacturing products and housing inventory to achieve \alert{economy of scales}
      \item[] \hskip3ex Number, Location, Capacity (manufacturing \& warehousing methodology)
    \end{itemize}
    }

\end{frame}

%=================================================================================================
\begin{frame}
  \frametitle{Example of decisions in Supply Chain}

  \begin{center}\small{
  \begin{tabular}{rlll} \hline \hline
    Area                      & \multicolumn{3}{c}{ \textbf{Levels} } \\ \cline{2-4}
    ~                         & Strategic         & Tactical          & Operation \\ \hline \hline
    Facility$^\pounds$           & number, size,     &                   &            \\
                              & location          &                   &  \\
    Inventory$^\pounds$          & form, policy      & safety stock      & qty, freq \\
                              &                   &                   &  \\
    Transportation$^\pounds$     & mode              & equipments        & routing \\
                              &                   &                   &  \\[6pt]
    Information$^\dag$        & OS, vendor        & interface         & process \\
                              &                   &                   &  \\
    Sourcing$^\dag$           & relationship~~    & selection, eval~~ & purchasing\\
                              &                   & contracts         &  \\
    Pricing$^\dag$            & scale, M/C        & mkt position      &  discount \\
                              & manuf. method     &                   &  \\    \hline \hline
  \end{tabular}
  }
  \end{center}

  \begin{flushleft}
    \begin{itemize}
      \item[] \code{$^\pounds$} \footnotesize{ \textbf{logistic driver:} physical movement $\multimap^{direct}\rightarrow$ operation }
      \item[] \code{$^\dag$} \footnotesize{ \textbf{cross-functional driver:} indirect related to physical movement  }
    \end{itemize}
  \end{flushleft}

  %%%------
  \note{
    \textbf{Purpose:} How to simplify decision making in SCM (cont)\\
    \textbf{Leading Question:} $<<$none$>>$
    \begin{itemize}
      \item \textbf{Inventory:} raw materials, WIP, F/G to \alert{smooth operations}
      \item[] \hskip3ex Cycle inventory  (expected demands); Safety inventory (unexpected demands); Seasonal inventory (predicted increasing demands) \NextHalfLine
      \item \textbf{Transportation:} moving products from $A$ to $B$
      \item[] \hskip3exMode of transportation; Route and network selection ; In-house \textbf{VS} outsource     \NextHalfLine
      \item \textbf{Information:} data and analysis to coordinate between stages (best drivers)
      \item[] \hskip3ex Push (MRP) \textbf{VS} pull (JIT); Coordination and information sharing;  Forecasting and aggregate planning ; Enabling technologies; IT infrastructure, KPI, assessing, evaluating parties \NextHalfLine
      \item \textbf{Sourcing:} functions a firm performs and functions that are outsourced
        \item[] \hskip3ex In-house \textbf{VS} outsource decisions ; Supplier evaluation and selection ; Procurement process ; identifying, certifying, and evaluating supplier of goods/services\NextHalfLine
      \item \textbf{Pricing:} Price associated with goods and services provided by a firm to the supply chain  $<<$unlikly topic for IE$>>$
      \item[] \hskip3ex pricing and qty ; Everyday low pricing \textbf{VS} high-low pricing ; fixed price \textbf{VS} menu pricing
    \end{itemize}
  }
\end{frame}

%=================================================================================================
\section[7-Eleven]{Case Study: 7-Eleven (Japan) Supply Chain}
%=================================================================================================
\begin{frame}
  \frametitle{7-Eleven Background}

    \begin{itemize}
      \item \textbf{Nature of business:} food and grocery Retailer
      \item \textbf{Status:} largest convenience store chain in Japan
      \item \textbf{Owned by:} Ito-Yokado Group (IYG) $\rightarrow$ Ito-Yokado (superstores) + Seven-Eleven Japan
      \item \textbf{Ex-Partner:} Southland Corp
      \item \textbf{Growths:} 1985 $\rightarrow$ 2007
      \begin{itemize}
        \item (519 in 1979) 2,299 $\rightarrow$ 12,304 stores
        \item 386 $\rightarrow$ 2,574 Billion JPY of Sales (7\% of IYG revenue)
        \item 9 $\rightarrow$ 91.5 Billion JPY of Income (47\% of IYG income)
        \item 12.6\% yearly avg revenue growth
        \item 20.9\%  yearly net income \NextHalfLine
      \end{itemize}
      \item \textbf{Performance:} decent even during recession (1990s) (why?)
    \end{itemize}



  \note{
    \textbf{Purpose:} 7Eleven JPN B/G and success story\\
    \textbf{Leading Question:} Why 7Eleven special in eyes of westerner \\
    \begin{itemize}
      \item \textbf{why 7Elven?} Ito-Yokado group also owns super stores. Guess what IYG needs from southland? (CP fresh mart \textsc{VS} 7Eleven)
      \item \textbf{history:} story of southland and impact to 7Eleven JPN
      \item \textbf{key idea of clip:} growth from acquiring and driving
      \item \textbf{Interesting Stat:}
      \begin{itemize} \scriptsize{
        \item 65\% of customer are male
        \item 55\% of customer live less than 5.0 minute
        \item 63\% of customer visit 2.0+ a week }
      \end{itemize}
      \item \textbf{lifestyle:} convenince store and JPN lifestyle
      \item \textbf{performance:} part timer who buys minimum only for survive after work. no car, no space, no time
    \end{itemize}
  }
\end{frame}

%=================================================================================================
\begin{frame}
  \frametitle{Convenience Store Industry and 7-Eleven}

  \textbf{Convenience Store Industry}
  \begin{description}
    \item[1991] 1.2\% (1991) of retailer outlet and 2.2\% (1991) of annual sales
    \item[1990s] 19,600 $\rightarrow$ 42,000 of retailer outlets (one for 3,000)
    \item[2004] 10 store chain $\approx$ 90\%
  \end{description}

  \NextLine
  \textbf{7-Eleven}
  \begin{itemize}
    \item \textbf{Big player:} 51\% of convenience store outlets ($2^{nd}$ Lawson by Mitsubishi)
    \item \textbf{Big Sales:} 31.5\% of total sales (30\% per store higher)
    \item \textbf{High Growth:} accounting for 60\% of industry growth
    \item \textbf{Focus products:} 70\% of total sales is food
  \end{itemize}

  \note{
    \textbf{Purpose:} Overview of 7Elven and Convenience industry\\
    \textbf{Leading Question:} $<<$none$>>$ \\
    \begin{itemize}
      \item \textbf{key idea of clip:} b/g of connection, stocks, distribution network, break-bulk
      \item \textbf{culture} USA $\leftrightarrow$ JPN $\rightarrow$ THA
      \item \textbf{risk} face-off energy conservation (24 hours operating) by safe place from stalker and criminal
    \end{itemize}
  }
\end{frame}

%=================================================================================================
\begin{frame}
  \frametitle{7-Eleven Operation}


  \begin{itemize}
    \item \textbf{Franchise System:} attractive (1\% awarded)
    \begin{itemize}
      \item<2-> new store \alert{within cluster of 50-60} old stores $\rightarrow$ limited present
      \item<2-> IYG got 45\% of total profits (68\% total revenues)
      \item<2-> 3 million JPY (training, initial inventory) \NextHalfLine
    \end{itemize}
    \item \textbf{Contents:} food items, beverage, magazine, consumer products
    \begin{itemize}
      \item<3-> high sales per area (3,000 SKU)
      \item<3-> \alert{food items} $\approx$ 40\% of total sales
      \item<3-> private soft drink \NextHalfLine
    \end{itemize}
    \item \textbf{Bill Payment Service:} convenient place to shop
    \begin{itemize}
      \item<4-> 10 million bills in 1994
      \item<4-> pick-up internet order at 7-Eleven
      \item<4-> utilization of IT
    \end{itemize}
  \end{itemize}


  \note{
    \textbf{Purpose:} operation component of 7Eleven JPN\\
    \textbf{Leading Question:} what does 7Eleven reveal so far?\\
    \begin{itemize}
      \item \textbf{Franchise:} good location is hard to own; why not partner? 'Right store' $\in$ walk, size, cluster, money
      \item \textbf{Content:} single uses, product analysis, private brand, 24/7 need energy and worker
      \item \textbf{Service:} service mail, utility payment, micro bank, money exchange, WiFi
      \item \textbf{Note} 7Catalog and 7pay was out-of-date and hacked. They eventually suspended
    \end{itemize}
  }
\end{frame}
%

%=================================================================================================
\begin{frame}
  \frametitle{7-Eleven IT}


  \begin{center}
    \only<1 |handout:1>{ \includegraphics[width=0.5\linewidth]{\figureLocation{scm_decision/info7eleven.jpg}} }
    \only<2-|handout:0>{
    \begin{tabular}{cc}
      \includegraphics[width=0.30\linewidth]{\figureLocation{scm_decision/gpTerm7eleven.jpg}}
    &
      \includegraphics[width=0.30\linewidth]{\figureLocation{scm_decision/ticketBuyer.jpg}}
    \end{tabular}
    }
  \end{center}

  \begin{flushright}
    \tiny{\textbf{source:} Attaran, M. 2004 ``Exploring the relationship between information technology and business process reengineering''}\\
  \end{flushright}

  \begin{itemize}
    \item \textbf{Graphic Order Terminal:} portable item analysis and order device
    \item \textbf{Scanner Terminal:} enable drop-off orders by drivers
    \item \textbf{Store computer:} processing house, link all terminals
    \item \textbf{POS Terminal:} record sales
  \end{itemize}


  \note{
    \textbf{Purpose:} IT as substitution of inventory \\
    \textbf{Leading Question:} $<<$none$>>$\\

    \begin{tabular}{ccc}
      \includegraphics[width=0.3\linewidth]{\figureLocation{scm_decision/gpTerm7eleven.jpg}}
      &
      \includegraphics[width=0.3\linewidth]{\figureLocation{scm_decision/pos7eleven.jpg}}
      &
      \includegraphics[width=0.3\linewidth]{\figureLocation{scm_decision/ticketBuyer.jpg}}
    \end{tabular}

    \begin{itemize}
      \item first to introduce point-of-sale (POS) information
      \item computer link to HQ and vendors $\rightarrow$ largest ISDN
      \item optical fiber to transfer voice, data, graphic information
      \item Hardware $\rightarrow$ inventory visibility
    \end{itemize}
  }
\end{frame}

%=================================================================================================
\begin{frame}
  \frametitle{7-Eleven Distribution System}

  \begin{center}
    \begin{tabular}{cc}
      \includegraphics[width=0.35\linewidth]{\figureLocation{scm_decision/truck7eleven.jpg}}
      &
      \includegraphics[width=0.35\linewidth]{\figureLocation{scm_decision/network7eleven.jpg}} \\
    \end{tabular}
  \end{center}
  \begin{flushright}
    \tiny{\textbf{source:} Jeffrey Hays 2009 \hskip10em}\\
  \end{flushright}

  \begin{itemize}
    \item short cycle time (3 for rices, 2 for fresh food)
    \item customize delivery time
    \item cut-off-time for breakfast, lunch, dinner
    \item combined delivery (products \& store)
    \item "Transfleet" provides transportation service
  \end{itemize}

  \note{
    \textbf{Purpose:} logistic of 7Eleven JN\\
    \textbf{Leading Question:} how a dedicate and customized distribution network look like?\\
    \begin{itemize}
      \item \textbf{info} 72 dedicated plants for food items
      \item \textbf{imply} delivery time \& multiple shipment $\rightarrow$ no checking
      \item \textbf{outsource} transportation, DC operator, DC location
    \end{itemize}
  }
\end{frame}

%=================================================================================================
\begin{frame}
  \frametitle{7-Eleven Outlet Locations}

  \begin{center}
    \only<1 |handout:0>{ \includegraphics [height=0.75\textheight]{\figureLocation{scm_decision/SevenElevenMap.png}} }
    \only<2 |handout:1>{ \includegraphics [height=0.75\textheight]{\figureLocation{scm_decision/SevenElevenMap02.png}} }
  \end{center}
  \begin{flushright}
    \tiny{\textbf{source:} Seven\& i HLDFS. Co~ \href{https://www.7andi.com/en/ir/library/co\_financial/2022/convenience\_store}{\code{https://www.7andi.com/en/ir/library/co\_financial/2022/convenience\_store}} \hskip10em}\\
  \end{flushright}

  \note{
    \begin{itemize}
      \item \textbf{cluster} High concentration $\rightarrow$ cannibalising its own share
    \end{itemize}
  }
\end{frame}

%=================================================================================================
\begin{frame}
  \frametitle{Study Questions}

  \begin{itemize}
    \item What are natures of business and core competencies?
    \item What are other rapid response supply chain similar?
    \item How does locations, inventory, IT, transportation support 7-Eleven SC strategy?
    \item Why do Seven Eleven stores offer many services?
    \item What are similarities and differences between 7Eleven Japan and 7Eleven Thailand
    \item What are \alert{trade-off} found in 7Eleven (comparing with Wal-mart)?
    \item<2-> What are concurrent issues about 7Eleven Japan? \hfill \code{\#\# out-of-scope}
    \item<2-> PRO and CON \alert{Couche-Tard} take-over? \hfill \code{\#\# out-of-scope}
  \end{itemize}


  \note{
    \textbf{Purpose:} example of question in case study\\
    \textbf{Leading Question:} $<<$none$>>$\\
    \begin{itemize}
      \item \textbf{What Core:} convenience food retailer (full price with IT)
      \item \textbf{Other rapid:} fast fashion, package delivery
      \item \textbf{SC Strategy:} premire location with people, few qty but many SKUs using IT and clustering
      \item \textbf{Other services:} $<<$reserch questions$>>$
      \item \textbf{Diff from Thai:} location and size of store-DC, food choices, \#delivery
      \item \textbf{Trade-Off:} location \textsc{VS} info \textsc{VS} qty
      \item \textbf{Thinking Points:} how do wal-mart and trader-Jo work?
    \end{itemize}
  }
\end{frame}

%=================================================================================================
\begin{frame}
  \frametitle{How to do a business case study}

  \begin{itemize}
    \item[1] Examine \alert{business environment} and \alert{information}
    \item[2] Analyze structure and size of \alert{main business}
    \item[3] Identify key issue or problem
    \item[4] Trace a chronological progression of how business responds
    \item[5] Learn successful aspects of responses as well as failures
    \item[6] \alert{Suggest alternative} successes, failures, unforeseen results, and measures
    \item[7] Describe \alert{what changes you would make} in business
  \end{itemize}

  \begin{flushright}
    \tiny{\textbf{source:} How to Analyse a Case Study (\href{http://www.wikihow.com}{http://www.wikihow.com}) \hskip10em}\\
  \end{flushright}


  \note{
    \textbf{Purpose:} guideline for case study \\
    \textbf{Leading Question:} how to do the case study?\\
    \begin{itemize}
      \item $<<$explain bullets$>>$
    \end{itemize}
  }
\end{frame}



%=================================================================================================
\section[Conceptual design]{Conceptual Designing of Distribution Network}
%=================================================================================================
\begin{frame}
  \frametitle{Role of Distribution in Supply Chain}


  \begin{center}
      \includegraphics[width=0.35\linewidth]{\figureLocation{scm_decision/distNetwork.jpg}}
  \end{center}


  \begin{itemize}
    \item Distribution $\rightarrow$ cost and customer experience; hence drives profitability
    \item Choice of distribution network $\rightarrow$ SC performance
    \item[] \hskip10em low cost $\leftrightarrow$ high responsiveness
  \end{itemize}

  \NextLine
  \textbf{Factors Influencing Network Design}
  \begin{itemize}
    \item Customer service
    \item Costs provided at level of service
  \end{itemize}

  \note{
    \textbf{Purpose:} trade-off between lv. customer service and cost \\
    \textbf{Leading Question:} After see example of 7Eleven distribution network, how do we apply the concept?
    \begin{itemize}
      \item \textbf{Role:} distribution network is strategic issue $\rightarrow$ daily shipment \& outsoucing
      \item \textbf{Selection of network:} cost, performance, operation, availability
      \item \textbf{What are design/thinking factors?}
    \end{itemize}
  }
\end{frame}

%=================================================================================================
\begin{frame}
  \frametitle{Trade-off: Areas of influencing}

  \begin{center}
      \begin{tabular}{ll}\hline
        Customer Service      & Cost \\ \hline
        \alert{Response time} & \alert{Inventories} \\
        Product variety       & \alert{Facilities \& handling cost}\\
        Product availability  & \alert{Transportation} \\
        Customer experience   & Information\\
        Order visibility      & \\
        Returnability         & \\ \hline
      \end{tabular}
  \end{center}


  \note{
    \textbf{Purpose:} service as qualitative \textsc{VS} cost as quantitative \\
    \textbf{Leading Question:} What are design/thinking factors?
    \begin{itemize}
      \item \textbf{Service:} with exception of \alert{time}, other is difficult to measure and vary
      \item \textbf{Cost:} money paid easy to measure
      \begin{itemize} \scriptsize{
        \item \textbf{facility:} rent, utility, training, admin, and infrastructure
        \item \textbf{inventory:} stock cost and opportunity
        \item \textbf{transportation:} movement and distribution network
        \item \textbf{information:} know cost but may or may not relate to \#stores as it depends on product and nature }
      \end{itemize}
    \end{itemize}
  }
\end{frame}

%=================================================================================================
\begin{frame}
  \frametitle{\# Facility \textbf{VS} Responding Time}

  \RemoveLine
  \begin{center}
    \includegraphics[width=\FigRatio\linewidth]{\figureLocation{scm_decision/Facility_RespTime.png}}\\
  \end{center}

  \begin{flushright}
    \tiny{\textbf{source:} Chopra and Meindl. 2010.~\cite{chopra2007supply} \hskip10em}\\
  \end{flushright}


  \note{
    \textbf{Purpose:} trade-off between \#store and response \\
    \textbf{Leading Question:} $<<$none$>>$\\
    \begin{itemize}
      \item \textbf{example} 400 Borders vs. 5 Amazon.com storage centers $\rightarrow$ wholefood and Amazon GO
    \end{itemize}
  }
\end{frame}

%=================================================================================================
\renewcommand{\FigRatio}{0.85}
%=================================================================================================
\begin{frame}
  \frametitle{Cost as function of \# facility}

  \RemoveLine

  \only<1|handout:0>{
%    \textbf{\# Facility \textbf{VS} Facility Cost  }
%    \RemoveLine
    \begin{center}
      \includegraphics[width=\FigRatio\linewidth]{\figureLocation{scm_decision/Facility_Facility.png}}
    \end{center}
  }

  \only<2|handout:0>{
%    \textbf{\# Facility \textbf{VS} Inventory Cost }
%    \RemoveLine
    \begin{center}
      \includegraphics[width=\FigRatio\linewidth]{\figureLocation{scm_decision/Invent_Facility.png}}
    \end{center}
  }

  \only<3|handout:0>{
%    \textbf{\# Facility \textbf{VS} Transportation Cost}
%    \RemoveLine
    \begin{center}
      \includegraphics[width=\FigRatio\linewidth]{\figureLocation{scm_decision/Trans_Facility.png}}
    \end{center}
  }

  \only<4-|handout:1>{
%    \textbf{\# Facility \textbf{VS} Total Costs}
%    \RemoveLine
    \begin{center}
    \includegraphics[width=\FigRatio\linewidth]{\figureLocation{scm_decision/Total_Facility.png}}
    \end{center}
  }

  \begin{flushright}
    \tiny{\textbf{source:} Chopra and Meindl. 2010.~\cite{chopra2007supply}\hskip10em}\\
  \end{flushright}

  \note{
    \textbf{Purpose:} Outline topics\\
    \textbf{Leading Question:} $<<$none$>>$\\
    \begin{itemize}
      \item \textbf{facility} Increase number of facilities $\rightarrow$ improve responsiveness.
      \item[] \hskip3ex If a firm is focusing on improving responsiveness, it may be beneficial for the company to increase the number of facilities beyond the total minimum cost point. Because better responsiveness $\rightarrow$ more revenue.
      \item \textbf{inventory:} reduction of safety stock from pooling effect
      \item \textbf{transportation} No distribution network will outperform others along all dimensions. Thus, it is important to ensure that the strengths of the distribution network fit with the strategic position of the firm
    \end{itemize}
  }
\end{frame}

%================================================================================================-
% Standardize how tikz displays vendors, crossdock, and stores
\usetikzlibrary{shapes}
\usetikzlibrary{arrows,positioning,decorations.pathmorphing}
\tikzstyle{factory}= [rectangle,draw=black!80,thick,fill=green!25,inner sep=0pt,minimum size=5mm]
\tikzstyle{retail}=  [diamond,draw=black!80,thick,fill=red!25,inner sep=0pt,minimum size=5mm,below]
\tikzstyle{customer}=[circle,draw=black!80,thick,fill=blue!25,inner sep=0pt,minimum size=5mm]
\tikzstyle{facility}=[rectangle,draw=black!80,thick,fill=yellow!25,text centered,text width=15mm,text height=3mm]
\tikzstyle{pFlow}=[->, black]
\tikzstyle{cFlow}=[draw=orange,ultra thick,decorate,decoration=snake]
\tikzstyle{dFlow}=[->,thick, blue, dotted]
\newcommand{\tikztext}[1]{\footnotesize{#1}}
%============================================================================================================--
\section[Distribution]{Conceptual Distribution Design }
%============================================================================================================--
\begin{frame}
  \frametitle{Distribution Design}

  \begin{center} \small{
  \begin{block}{Question: How to buy and pay for these items? }{}
    \begin{itemize}
      \item \textbf{FMCG:}: shampoo, shaving razor \hfill \uncover<2->{\code{anywhere}}
      \item \textbf{Hi-End Consumer:}: iPhone, lenovo X1 \hfill \uncover<2->{\code{few place}}
      \item \textbf{Luxury Products:} ferrari, home appliance, jewelry
      \item \textbf{One-Of-Kinds:} machine, house, event
    \end{itemize}
  \end{block} }
  \end{center}


  \begin{itemize}
    \item \textbf{Channel of Distribution:} network of firms collaborates to provide products to consumers
    \item \textbf{Channel type:}
    \begin{itemize}
      \item \textbf{distributor:} \alert{firms} who provides assortment, break-bulk, finance, and inventory
      \item \textbf{value added retailer (VAR):} \alert{specialists} who works with manufacturers
      \item \textbf{broker/agent:} \alert{individuals} who represent manufacturers  \NextHalfLine
    \end{itemize}
    \item \textbf{Network:} intensive (many-in-TH) $\gg$ selective $\gg$ exclusive (1-in-TH)
  \end{itemize}

  \note{
    \textbf{Purpose:} Overview of how to distribute\\
    \textbf{Leading Question:} how char. of product affect distribution network and its development?\\
    \begin{itemize}
      \item \textbf{reverse phycology:} rarity and price as signal of class
      \item \textbf{Type of Mkt distribution Network:}
      \begin{itemize} \scriptsize{
        \item \textbf{intensive:} available and \alert{saturated} everywhere; otherwise lost sale. e.g., common FMCG
        \item \textbf{selected:} available in some regions or some store. partner with training e.g., cosmetic
        \item \textbf{extensive:} extreme case of selective. e.g. [ex] hermes, rolex, super car     }
      \end{itemize}
      \item \textbf{note:} how to track information flow in supply chain
      \item \textbf{more info} \href{https://www.yourarticlelibrary.com/distribution/types-of-distribution-intensive-selective-and-exclusive-distribution/5780}{\code{more info}}
    \end{itemize}
  }
\end{frame}

%============================================================================================================--
\begin{frame}
  \frametitle{Direct shipments}

  \begin{center}
    \begin{tikzpicture}[xscale=4,yscale=2.5]

      \node[factory] (f1) at (0,1.0) {};
      \node[factory] (fn) at (0,0.0) {};
      \node[customer] (c1) at (2.5,1.0) {};
      \node[customer] (c2) at (2.5,0.5) {};
      \node[customer] (cn) at (2.5,0.0) {};

      %\node at (0,0.5) {$\vdots$};
      \node at (0,-0.5) {\tikztext{Factories}};

      %\node at (1,0.5) {$\vdots$};
      \node at (2.5,-0.5) {\tikztext{Customer}};

      \draw [pFlow] (f1) -- (c1);
      \draw [pFlow] (f1) -- (c2);
      \draw [pFlow] (f1) -- (cn);
      \draw [pFlow] (fn) -- (c1);
      \draw [pFlow] (fn) -- (c2);
      \draw [pFlow] (fn) -- (cn);

      \node[retail]  (r1) at (1.0,1.55) {};
      \node[above=0.2cm of r1] (rName) {\tikztext{Retailer}};
      \draw [dFlow] (c1) -- (r1);
      \draw [dFlow] (c2) -- (r1);
      \draw [dFlow] (cn) -- (r1);
      \draw [dFlow] (r1) -- (f1);
      \draw [dFlow] (r1) -- (fn);

      \draw[-,dotted](0.5,-0.35) -- node[pos=0.5,auto] {legend} ++(1.05,0.00) -- ++(0.0,-0.3) -- ++(-1.05,0.0) -- ++(0.0,0.3) ;
      \draw[pFlow] (0.6,-0.55) -- node[pos=0.5,auto] {\tikztext{product}} ++(0.25,0.0);
      \draw[dFlow] (1.1,-0.55) -- node[blue, pos=0.5,auto] {\tikztext{info}} ++(0.25,0.0);

    \end{tikzpicture}
  \end{center}

  \note{
    \textbf{Purpose:} show direct shipment\\
    \textbf{Leading Question:} $<<$none$>>$\\
    \begin{itemize}
      \item explain flow of phy and info
      \item simple
    \end{itemize}
  }
\end{frame}


%================================================================================================-
\begin{frame}
  \frametitle{Direct shipping with in-transit merge}


  \begin{center}
    \begin{tikzpicture}[xscale=4,yscale=2.5]

      \node[factory] (f1) at (0,1.0) {};
      \node[factory] (fn) at (0,0.0) {};
      \node[facility] (facility) at (1.75,0.5) {\tikztext{Transit}};
      \node[customer] (c1) at (2.5,1.0) {};
      \node[customer] (c2) at (2.5,0.5) {};
      \node[customer] (cn) at (2.5,0.0) {};

      %\node at (0,0.5) {$\vdots$};
      \node at (0,-0.5) {\tikztext{Factories}};

      %\node at (1,0.5) {$\vdots$};
      \node at (2.5,-0.5) {\tikztext{Customer}};

      \draw [pFlow] (f1) -- (facility);
      \draw [pFlow] (fn) -- (facility);
      \draw [pFlow] (facility) -- (c1);
      \draw [pFlow] (facility) -- (c2);
      \draw [pFlow] (facility) -- (cn);

      \node[retail]  (r1) at (1.0,1.55) {};
      \node[above=0.2cm of r1] (rName) {\tikztext{Retailer}};
      \draw [dFlow] (c1) -- (r1);
      \draw [dFlow] (c2) -- (r1);
      \draw [dFlow] (cn) -- (r1);
      \draw [dFlow] (r1) -- (f1);
      \draw [dFlow] (r1) -- (fn);

      \draw[-,dotted](0.5,-0.35) -- node[pos=0.5,auto] {legend} ++(1.05,0.00) -- ++(0.0,-0.3) -- ++(-1.05,0.0) -- ++(0.0,0.3) ;
      \draw[pFlow] (0.6,-0.55) -- node[pos=0.5,auto] {\tikztext{product}} ++(0.25,0.0);
      \draw[dFlow] (1.1,-0.55) -- node[blue, pos=0.5,auto] {\tikztext{info}} ++(0.25,0.0);

    \end{tikzpicture}
  \end{center}

  \note{
    \textbf{Purpose:} how to modify direct for response \\
    \textbf{Leading Question:} <<none>>\\
    \begin{itemize}
      \item explain flow of phy and info
      \item Both slide are a trading company with/without hub
      \item retail needs to coordinate shipment
      \item retail does not own transit facility
      \item increasing effort
    \end{itemize}
  }
\end{frame}

%================================================================================================-
\begin{frame}
  \frametitle{Carrier delivery by distributors}


  \begin{center}
    \begin{tikzpicture}[xscale=4,yscale=2.5]

      \node[factory] (f1) at (0,1.0) {};
      \node[factory] (fn) at (0,0.0) {};
      \node[facility] (facility) at (1.25,0.5) {\tikztext{Warehouse}};
      \node[customer] (c1) at (2.5,1.0) {};
      \node[customer] (c2) at (2.5,0.5) {};
      \node[customer] (cn) at (2.5,0.0) {};

      %\node at (0,0.5) {$\vdots$};
      \node at (0,-0.5) {\tikztext{Factories}};

      %\node at (1,0.5) {$\vdots$};
      \node at (2.5,-0.5) {\tikztext{Customer}};

      \draw [pFlow] (f1) -- (facility);
      \draw [pFlow] (fn) -- (facility);
      \path [pFlow] (facility.east) edge[bend left] (c1);
      \path [pFlow] (facility.east) edge[bend left] (c2);
      \path [pFlow] (facility.east) edge[bend left] (cn);

      %\node[retail]  (r1) at (1.0,1.55) {};
      %\node[above=0.2cm of r1] (rName) {\tikztext{Retailer}};
      \path[dFlow] (c1) edge[bend left] (facility.east);
      \path[dFlow] (c2) edge[bend left] (facility.east);
      \path[dFlow] (cn) edge[bend left] (facility.east);


      \draw[-,dotted](0.5,-0.35) -- node[pos=0.5,auto] {legend} ++(1.05,0.00) -- ++(0.0,-0.3) -- ++(-1.05,0.0) -- ++(0.0,0.3) ;
      \draw[pFlow] (0.6,-0.55) -- node[pos=0.5,auto] {\tikztext{product}} ++(0.25,0.0);
      \draw[dFlow] (1.1,-0.55) -- node[blue, pos=0.5,auto] {\tikztext{info}} ++(0.25,0.0);

    \end{tikzpicture}
  \end{center}

  \note{
    \textbf{Purpose:} further modification with inventory \\
    \textbf{Leading Question:} <<none>>\\
    \alert{Intro}
    \begin{itemize}
      \item there is inventory
    \end{itemize}
  }
\end{frame}

%================================================================================================-
\begin{frame}
  \frametitle{Last mile delivery by distributors}


  \begin{center}
    \begin{tikzpicture}[xscale=4,yscale=2.5]

      \node[factory] (f1) at (0,1.0) {};
      \node[factory] (fn) at (0,0.0) {};
      \node[facility] (facility) at (1.25,0.5) {\tikztext{Warehouse}};
      \node[customer] (c1) at (2.5,1.0) {};
      \node[customer] (c2) at (2.5,0.5) {};
      \node[customer] (cn) at (2.5,0.0) {};

      %\node at (0,0.5) {$\vdots$};
      \node at (0,-0.5) {\tikztext{Factories}};

      %\node at (1,0.5) {$\vdots$};
      \node at (2.5,-0.5) {\tikztext{Customer}};

      \draw [pFlow] (f1) -- (facility);
      \draw [pFlow] (fn) -- (facility);
      \draw [pFlow,red] (c1) -- (c2);
      %\draw [pFlow] (c2) -- (cn);
      \path [pFlow,red] (facility.east) edge[bend left] (c1);
      \path [pFlow,red] (c2.south) edge[bend left] (facility.east);
      \path [pFlow,red] (facility.east) edge[bend left] (cn.north);
      \path [pFlow,red] (cn.south) edge[bend left]  (facility.east);

      %\node[retail]  (r1) at (1.0,1.55) {};
      %\node[above=0.2cm of r1] (rName) {\tikztext{Retailer}};
      \path[dFlow] (c1) edge[bend left] (facility.east);
      \path[dFlow] (c2) edge[bend left] (facility.east);
      \path[dFlow] (cn) edge[bend left] (facility.east);


      \draw[-,dotted](0.5,-0.35) -- node[pos=0.5,auto] {legend} ++(1.05,0.00) -- ++(0.0,-0.3) -- ++(-1.05,0.0) -- ++(0.0,0.3) ;
      \draw[pFlow] (0.6,-0.55) -- node[pos=0.5,auto] {\tikztext{product}} ++(0.25,0.0);
      \draw[dFlow] (1.1,-0.55) -- node[blue, pos=0.5,auto] {\tikztext{info}} ++(0.25,0.0);

    \end{tikzpicture}
  \end{center}

  \note{
    \textbf{Purpose:} further modification with transportation \\
    \textbf{Leading Question:} <<none>>\\
    \begin{itemize}
      \item improve transportation
    \end{itemize}
  }
\end{frame}

%================================================================================================-
\begin{frame}
  \frametitle{Customer pick up}


  \begin{center}
    \begin{tikzpicture}[xscale=4,yscale=2.5]

      \node[factory] (f1) at (0,1.0) {};
      \node[factory] (fn) at (0,0.0) {};
      \node[facility] (p1) at (1.85,0.75) {\tikztext{site \#1}};
      \node[facility] (pn) at (1.85,0.25) {\tikztext{site \#2}};
      \node[facility] (dock) at (1.0,0.75) {\tikztext{Crossdock}};
      \node[customer] (c1) at (2.5,1.0) {};
      \node[customer] (c2) at (2.5,0.5) {};
      \node[customer] (cn) at (2.5,0.0) {};

      %\node at (0,0.5) {$\vdots$};
      \node at (0,-0.5) {\tikztext{Factories}};

      %\node at (1,0.5) {$\vdots$};
      \node at (2.5,-0.5) {\tikztext{Customer}};

      \uncover<2->{
      \draw [pFlow] (f1) -- (dock.west);
      \draw [pFlow] (fn) -- (dock.west);
      \draw [pFlow] (dock.east) -- (p1.west);
      \draw [pFlow] (dock.east) -- (pn.west);
      \path [cFlow] (c1.west) -- (p1);
      \path [cFlow] (c2.west) -- (pn);
      \path [cFlow] (cn.west) -- (pn);
      }

      \node[retail]  (r1) at (1.0,0.25) {};
      \node at (1.0,0.05) {\tikztext{Retailer}};
      \draw[dFlow] (c1) -- (r1.east);
      \draw[dFlow] (c2) -- (r1.east);
      \draw[dFlow] (cn) -- (r1.east);
      \draw[dFlow] (r1.west) -- (f1);
      \draw[dFlow] (r1.west) -- (fn);



      \draw[-,dotted](0.5,-0.55) -- node[pos=0.5,auto] {legend} ++(1.05,0.00) -- ++(0.0,-0.3) -- ++(-1.05,0.0) -- ++(0.0,0.3) ;
      \draw[pFlow] (0.6,-0.75) -- node[pos=0.5,auto] {\tikztext{product}} ++(0.25,0.0);
      \draw[dFlow] (1.1,-0.75) -- node[blue, pos=0.5,auto] {\tikztext{info}} ++(0.25,0.0);

    \end{tikzpicture}
  \end{center}

  \note{
    \textbf{Purpose:} current practice and origin of omni-channel \\
    \textbf{Leading Question:} $<<$none$>>$\\
    \begin{itemize}
      \item increase flexibility, like 7Eleven JPN
    \end{itemize}
  }
\end{frame}



%================================================================================================-

\begin{frame}
  \frametitle{Reference}

  \begin{footnotesize}
    \bibliography{\refLocation{Coordination}}
  \end{footnotesize}

  \note{
    [empty] note for citation page
  }

\end{frame}

%=================================================================================================
\end{document}









%=================================================================================================
\begin{frame}
  \frametitle{Frame Title Here}

  \note{
    \begin{itemize}
      \item info
    \end{itemize}
  }
\end{frame}

%=================================================================================================
\begin{frame}
  \frametitle{Responsibility}

    \begin{center}
    \begin{tabular}{rr} \hline
      7-Eleven   & Owner   \\ \hline
      \alert{Supply products}
    &
      \alert{Operate} \& \alert{maintain} store
    \\
      Accounting, Advertising, IT Systems
    &
      Hiring staff \& customer service
    \\
      Install and remodel
    &
      Order products
    \\ \hline
    \end{tabular}
    \end{center}

  \note{
    \begin{itemize}
      \item info
    \end{itemize}
  }
\end{frame}

%=================================================================================================
\begin{frame}
  \frametitle{Though exercises}

  Reason for outsoucing
 Restructuring the supply bass
•  Capital circumstance no longer the
same
•  Shift on “what we are good at”
•  Change to world capabilities
•  More pressure on cost control
•  Resources no-longer available internally
•  Customer demand changed


  \begin{alertblock}{}{}
    Name a company that \alert{bad} in supply chain and elaborate
  \end{alertblock}

  \NextHalfLine
  \textsc{What are characteristic of \alert{bad} supply chain?}
  \begin{itemize}
    \item \textbf{Product:} items, quantities, condition
    \item \textbf{Time:} delay, early
    \item \textbf{Cost:}
    \item \textbf{Place:}
  \end{itemize}


  \note{
    \begin{itemize}
      \item important and motivation of SCM~\cite{jacobs2010operations}~\cite{ballou2007evolution} \cite{chopra2007supply} \cite{goetschalckx2011supply}
      \item SCM is an elective class/key application of ICE
    \end{itemize}
  }
\end{frame}

%=================================================================================================
\begin{frame}
  \frametitle{Factors affected supply chain}

  \begin{itemize}
    \item \textbf{Supply management:} \alert{forecasting}, sourcing, supplier evaluation, supplier certification, strategic partnership
    \item \textbf{Production management:} \alert{production planning} (flow/job shop), \alert{JIT \textsc{VS} MRP}, \alert{Inventory planning}
    \item \textbf{Transportation management:} \alert{mode selection}, \alert{distribution network}
    \item \textbf{Integration:} Process integration
  \end{itemize}

  \NextLine
  \begin{alertblock}{Message}{}
    Every activities affect efficiency and responsiveness of supply chain
  \end{alertblock}

  \note{
    \begin{itemize}
     \item \alert{Production planning:} Aggregated planning, S\&OP
     \item \alert{Inventory management}: EOQ (deterministic), newsvendor (perishable \& stochastic)
    \end{itemize}

  \textbf{Trends}
  \begin{itemize}
    \item Supplier improvement, e.g., evaluation and certification
    \item Long-term relationship and information sharing
    \item Strategic partnerships = successful and trusting relationships with top-performing suppliers
    \item In-house or outsourcing
    \item Flexible Vs Dedicated
  \end{itemize}

  \begin{itemize}
    \item \textbf{Coordinating between separated organizations:} wake of outsourcing $\rightarrow$ complex supply chain. Hence, better communication in product process development, manufacturing, packaging, testing, after sale services \NextHalfLine
    \item \textbf{Optimizing Global Supply Chain:} expanding partnerships and building facilities in foreign markets.
    \begin{itemize}
      \item \textbf{breadth} foreign manufacturing, office \& retail sites, foreign suppliers \& customers
      \item \textbf{depth} involving second and third tier suppliers \& customers \NextHalfLine
    \end{itemize}
    \item \textbf{Supply Chain Responsiveness:} more flexible and responsive to customer needs; required industry benchmark and continuous improve \NextHalfLine
    \item \textbf{Social Responsibility:} creating \alert{long-lasting brand} by \alert{triple bottom line}
    \begin{itemize}
      \item \textbf{economic} profit, value
      \item \textbf{employee} job security, positive work condition, equal opportunity
      \item \textbf{environmental} carbon emission, recycle, green logistic
    \end{itemize}
  \end{itemize}

  }
\end{frame} 